\chapter{Enhanced Player Play}
łabel{ch:enhanced−play}

The world of \textbf{Fate's Edge} sings when players help steer the tune. This chapter gathers \emph{player−facing dials}−−−lightweight, opt‑in mechanics that turn your table's choices into momentum. Treat these as a menu: use a few, or layer many as your group grows comfortable.

\section{Role‑Based Tactics}
łabel{sec:role−tactics}
∈dex{tactics!role−based}

Your role in the party determines how to best apply pressure, conserve resources, and turn the tide. Use these tactical guidelines regardless of your specific build.

\subsection{Melee Striker (Duelist / Assassin / Berserker)}
\textit{Your job is to eliminate high‑value targets and break enemy lines.}
\begin{itemize}
  ℹtem \textbf{Be aggressive.} You move faster than magic and hit harder than ranged fire. Close the distance early; every round you spend not engaged is a round you are not doing your job.
  ℹtem \textbf{Flank, flank, flank.} A flanked enemy loses Position advantages. Coordinate with another melee character to attack from two sides.
  ℹtem \textbf{Use Boons on re‑rolling failed attacks, not defense.} Your HP is a resource; their death is the objective.
  ℹtem \textbf{When outnumbered, use the terrain.} A narrow corridor or doorway lets you fight one enemy at a time while allies shoot over your shoulder.
  ℹtem \textbf{Know your exit.} If you go down, the party loses its hammer. Keep a backup plan (a smoke bomb, a doorway to retreat through, a Bonded ally who can pull you out).
\end{itemize}

\subsection{Defender (Guardian / Shield / Tank)}
\textit{You hold the line so others don't have to.}
\begin{itemize}
  ℹtem \textbf{Use the Defend action early.} It improves your Position by one step and gives you a free reaction to intercept attacks on allies.
  ℹtem \textbf{Spend Boons on re‑rolling defense failures.} Your survival keeps the party safe; a missed attack on your part is less catastrophic than a missed defense.
  ℹtem \textbf{Position yourself between the enemy and your squishy allies.} The GM will honour the fiction −− a foe cannot shoot past you without moving around you.
  ℹtem \textbf{Trade Harm for Fatigue whenever possible.} A wound that clears your Fatigue track is often worth it.
  ℹtem \textbf{Protect at least one ally with a Bond.} That Bond gives you a free Boon refund when you assist them −− perfect for \emph{Cover & Protect}.
\end{itemize}

\subsection{Controller (Mage / Crowd Control / Terrain Shaper)}
\textit{You shape the battlefield, forcing enemies into bad positions.}
\begin{itemize}
  ℹtem \textbf{Open with area control.} A \texttt{[AREA]} \texttt{[BIND]} or \texttt{[WALL]} spell can split the enemy force, letting your strikers pick them off.
  ℹtem \textbf{Use Boons to improve Position before casting, not after.} A Controlled cast is much safer than a Desperate one.
  ℹtem \textbf{Accept SB as a resource, not a penalty.} You generate them anyway −− lean into it. Offer the GM a Devil's Bargain to make your spell more effective in exchange for a later complication.
  ℹtem \textbf{Tags matter.} \texttt{[EARTH]} + \texttt{[WALL]} blocks movement; \texttt{[AIR]} + \texttt{[LEAP]} repositions allies; \texttt{[FATE]} + \texttt{[OMEN]} gives the party a clue. Read the Tag reference (\ref{ch:magic}).
  ℹtem \textbf{Do not waste spells on enemies your strikers can kill in one hit.} Save your power for clusters, bosses, or emergencies.
\end{itemize}

\subsection{Support (Healer / Buffer / Bard)}
\textit{You multiply the party's effectiveness.}
\begin{itemize}
  ℹtem \textbf{First action: improve Position or grant a Boon to a striker.} A Dominant striker kills faster than you can heal.
  ℹtem \textbf{Use \emph{Inspire} early.} The +1 die and extra Boon can turn a fight.
  ℹtem \textbf{Heal Harm after the fight, Fatigue during it.} \texttt{[HEAL]} as a cantrip clears Fatigue (safe and fast); healing Harm requires a full action and may generate SB.
  ℹtem \textbf{Carry a ranged weapon.} You do not need to be a damage dealer, but plinking a foe can distract them or force a reaction.
  ℹtem \textbf{Your Bonds are your fuel.} Transfer Boons to the striker who needs them. A bonded ally gets a Boon refund on their first assist −− use it.
\end{itemize}

\subsection{Scout (Rogue / Ranger / Infiltrator)}
\textit{You see the danger before it sees the party.}
\begin{itemize}
  ℹtem \textbf{Go first.} Before the fight, scout ahead. The information you bring is worth more than an extra sword.
  ℹtem \textbf{Strike from stealth, then vanish.} \emph{Backstab} + \emph{Shadow Dance} + \emph{Deathblow} is a classic chain (see \ref{ch:talents}).
  ℹtem \textbf{Use the environment.} High ground, cover, and shadow are your friends. Spend a Move action to reposition even if you do not attack.
  ℹtem \textbf{Your independent action (Scout & Signal) is free Position for the party.} Use it before the first roll of a fight.
  ℹtem \textbf{Do not be the first to engage in a fair fight.} If you are trading blows, something went wrong.
\end{itemize}

\subsection{Leader (Commander / Tactician / Face)}
\textit{You turn a group of individuals into a unit.}
\begin{itemize}
  ℹtem \textbf{Setup actions are your bread and butter.} Grant an ally +1 Position or step up Effect −− no roll, just fiction.
  ℹtem \textbf{Use \emph{Command} to rally.} A successful Command roll can remove Fear conditions or improve morale.
  ℹtem \textbf{Delegate.} Your followers (see \ref{ch:followers}) are not just dice −− they can take independent actions. Trust them.
  ℹtem \textbf{Call shots.} ``Focus the mage,'' ``Take cover,'' ``Fall back'' −− the GM will adjust Position based on coherent orders.
  ℹtem \textbf{Your Boons are best spent on assisting others.} You are a force multiplier. A Boon spent on a striker's attack is worth more than your own attack.
\end{itemize}

\section{Tactical Play for Players}
łabel{sec:tactical−play}

Combat and high‑stakes scenes reward clever thinking as much as raw dice. Use these tactics to turn the fiction in your favor.

\subsection{Closing on Ranged and Magic Users}

A bowman on a balcony or a mage behind a ward can shred your party from a distance. Close the gap without getting shredded.

\begin{itemize}
  ℹtem \textbf{Use Cover and Movement.} Move from cover to cover. Spend a \textbf{Move action} to dash between obstacles. Ending behind solid cover may let the GM set your next defense roll as \textbf{Controlled} instead of Desperate.
  ℹtem \textbf{Create Distractions.} Spend \textbf{1 Boon} to describe a diversion. On a successful \textbf{Deception} or \textbf{Command} roll, the ranged attacker may lose their aim or shift attention for one exchange, dropping your Position penalty by one step.
  ℹtem \textbf{Teamwork: The Running Shield.} Have a heavily armoured ally take the \textbf{Defend} action while moving forward. They intercept shots (converting Harm to Fatigue) while you sprint behind them. The GM may rule that the defender blocks line of sight, granting you \textbf{Dominant} position for your next approach.
  ℹtem \textbf{Suppress Them.} Use the \textbf{Suppress} option (core rules). You don't need to hit −− just force the enemy to keep their head down. The GM may impose a −−1 die penalty on their next attack or worsen their Position.
  ℹtem \textbf{Spend Boons to Close.} Spend \textbf{1 Boon} to improve your Position by one step before rolling movement or defense. That can mean the difference between a bolt in the chest and a glancing hit.
\end{itemize}

\paragraph{Example.}
Lyra (rogue) wants to cross a courtyard to reach a mage on a balcony. She spends 1 Boon to improve her Position from Desperate to Controlled, then stays close to a crumbling wall. The GM rules that the mage's magic missile is at −−1 die because of partial cover. Lyra makes it across with a Partial success, taking Fatigue instead of Harm.

\subsection{When You're Outmatched}

Sometimes the enemy is bigger, faster, or more numerous. You can still win −−− or survive.

\begin{itemize}
  ℹtem \textbf{Don't Trade Blows.} Instead of attacking, use your action to \textbf{Setup} for an ally (grant them +1 Position or step up Effect). Spend 1 Boon to assist them directly (+1 die). One well‑supported PC can turn a fight faster than two desperate ones.
  ℹtem \textbf{Accept Story Beats (SB) to Buy Time.} When you roll a Partial or Miss, the GM gets SB. You can proactively offer to accept \textbf{+1 SB} (banked, not spent now) in exchange for a short narrative delay: ``I fall back behind a pillar −− they lose sight of me for a moment.'' This doesn't change the roll's outcome, but gives you a breath to reposition.
  ℹtem \textbf{Spend Boons on Defense, Not Offense.} When outnumbered, a missed attack hurts less than a received hit. Save Boons to re‑roll failed defence rolls or to improve your Position before an enemy's big swing. A \textbf{Boon spent to re‑roll a 1 on a defence roll} can prevent Harm 2 or 3.
  ℹtem \textbf{Escape and Reset.} There's no shame in running. Use the \textbf{Withdraw} action (core rules). Spend 1 Boon to make it a \textbf{Controlled} withdrawal instead of Desperate. Live to fight another scene.
  ℹtem \textbf{Use the Environment as a Weapon.} Ask the GM: ``Is there a chandelier? A sand pile? A steep stair?'' Describe how you trigger it. The GM may set a low DV (2−−3); on success, the obstacle deals Harm, imposes a Condition, or improves the party's Position for a round.
\end{itemize}

\paragraph{Example.}
Benjamin (fighter) faces three bandits. Instead of attacking, he spends 1 Boon to assist the group's archer, then kicks over a brazier. The GM calls for a Body + Athletics roll (DV 2). Benjamin succeeds −− the fire creates a smoky screen, worsening the bandits' Position to Desperate for one exchange. The archer picks one off.

\subsection{When to Push It −−− and When Not To}

The \textbf{Push It} mechanic (usually on Rites or Cantor Songs) escalates an effect immediately at a cost: extra Obligation, Fatigue, Corruption, or a GM Story Beat. Use it wisely.

⫕section*{Good Times to Push}
\begin{itemize}
  ℹtem \textbf{Preventing a Party Wipe.} If a teammate is about to take lethal Harm, Push to get that heal or ward now.
  ℹtem \textbf{Ending a Scene.} The timer is about to fill, or the enemy leader is fleeing. Push to finish it.
  ℹtem \textbf{When the Cost Is Acceptable.} If you have spare Fatigue, or you're already at max Obligation and about to trigger a Patron demand anyway, the extra cost doesn't hurt as much.
\end{itemize}

⫕section*{Avoid Pushing When}
\begin{itemize}
  ℹtem \textbf{You're Already in Trouble.} Pushing adds Fatigue or SB. If you're already at Fatigue 2 or have a full Death Timer , the extra cost could break you.
  ℹtem \textbf{A Partial Success Is Enough.} Many Rites and Songs still work on a Partial −− just weaker or shorter. Don't pay extra for power you don't require.
  ℹtem \textbf{The GM Is Low on SB.} If the GM has 0−−1 SB banked, they can't immediately punish a miss. But if they have 3+, a failed Push could trigger a major twist. Check the table's SB tracker.
\end{itemize}

⫕section*{How to Avoid Needing to Push}
\begin{itemize}
  ℹtem \textbf{Use Boons Early.} Spend 1 Boon to improve your Position before the roll. A Controlled roll is less likely to fail than a Desperate one.
  ℹtem \textbf{Get Help.} Have an ally assist you (spend 1 Boon for +1 die). That +1 die might turn a Partial into a Clean Success.
  ℹtem \textbf{Lower the DV with Fiction.} ``I wait for the guards to change shift'' might lower DV by 1. ``I use my Scholar's Cell free use to research the ward's pattern'' could drop DV from 4 to 3.
  ℹtem \textbf{Accept Failure and Bank Boons.} A Miss gives you 2 Boons. Sometimes it's better to fail now, take the Boons, and try again next round with a bigger pool or better Position.
\end{itemize}

\paragraph{Example.}
Sariel (Cantor) needs to charm a guard quickly. She could Push to resolve the Song instantly, but she's already at Fatigue 2. Instead, she spends 1 Boon to improve her Position from Controlled to Dominant, then has an ally whisper a distraction (assist, another Boon). She rolls 4 dice (Lore + Performance) against DV 3. She gets 3 successes −− a Clean Success. No Push needed, and she saved her Fatigue.

\begin{tcolorbox}[colback=blue!5!white,colframe=blue!75!black,title=Reminder]
You can spend \textbf{1 Boon} before a roll to increase Effect by one step.
\end{tcolorbox}

\subsection{General Player Tips for Enhanced Experience}

\begin{itemize}
  ℹtem \textbf{Track Your Boons, Fatigue, and Obligation.} These are your resources. A character with 4 Boons and 0 Fatigue is much stronger than one with 0 Boons and Fatigue 2.
  ℹtem \textbf{Use Bonds to Shift Boons.} At Tier III, you can transfer up to 2 Boons to a bonded ally as a free action. Fuel the party's best character for a crucial roll. Don't hoard.
  ℹtem \textbf{Communicate Your Intent.} Tell the GM what you want to achieve, not just what skill you're rolling. ``I want to scare him into dropping the keys'' is better than ``I roll Intimidation.'' The GM can set a fairer DV and offer better consequences.
  ℹtem \textbf{Ask About the Environment.} ``Is there loose gravel? A tapestried wall? A rain barrel?'' These become levers. The GM wants you to use them.
  ℹtem \textbf{Know When to Fold.} Not every fight is winnable. Spending 2 Boons to escape with everyone alive is a victory. The game rewards survival and smart choices.
  ℹtem \textbf{Use Downtime.} Between sessions, update your logs, tend your Assets, and clear Complications. A well‑maintained Asset (Neglected → Maintained) can save your life next session. A cleared Complication removes a ``free SB'' button from the GM.
  ℹtem \textbf{Talk to Your GM.} If a rule feels off or a scene drags, speak up. Fate's Edge is collaborative. The GM can adjust DV, offer alternative stakes, or introduce a Devil's Bargain to keep things moving.
\end{itemize}

\section{Player Resources}
∈dex{resources!player}∈dex{Boons}∈dex{Story Beats}

Two shared currencies appear throughout this chapter:
\begin{description}
  ℹtem[\textbf{Boons}] Core system edge; spend to power talents or convert to XP (see \ref{ch:core−mechanics}).
  ℹtem[\textbf{SB}] The GM's Story Beats; several options here invite players to accept SB in exchange for narrative benefits.
\end{description}

\section{Engagement Rewards}
∈dex{engagement}∈dex{investment tracker}

\subsection{Session Investment Tracker}
∈dex{investment tracker}
At session close, each player privately rates their \emph{investment} (1−−3). The tracker rewards steady participation without penalising quiet nights.

\begin{description}
  ℹtem[\textbf{1 −−− Low}] You followed others' lead. Gain \textbf{+1d} once next session on any \emph{support} action.
  ℹtem[\textbf{2 −−− Medium}] You took initiative in some scenes. Gain \textbf{+1d} once next session on a \emph{relationship} roll and bank \textbf{1 Momentum} (see \ref{subsec:momentum−banking}).
  ℹtem[\textbf{3 −−− High}] You anchored or elevated multiple scenes. Gain \textbf{+1d} twice next session (different scenes) and bank \textbf{1 Momentum}.
\end{description}

\paragraph{Table Boon.}
If \emph{all} players report 2+ in a session, the GM may award the table \textbf{1 free Boon} to assign at the start of the next session.

\subsection{Cultural Immersion Bonus}
∈dex{culture!immersion}∈dex{roleplay rewards}
Reward lived‑in play with light, predictable benefits.
\begin{itemize}
  ℹtem \textbf{Earn 1 Culture Point} when you enrich a scene with apt language, rites, or customs (Sea‑Patter hail; bell‑speech courtesy; Aeler craft‑honorifics). Max \textbf{3/session}.
  ℹtem \textbf{Spend 3 Culture} for \textbf{+1d} on a culturally keyed roll.
  ℹtem \textbf{Spend 5 Culture} (across sessions) to declare a \emph{trusted door}: one institution in that culture treats you as \textit{Friendly} for the next approach.
\end{itemize}

\section{Collaborative Play}
∈dex{collaborative play}∈dex{negotiation}

\subsection{Information Trading}
∈dex{information economy}
When the table hunts answers, convert curiosity into structure.

\paragraph{Request & Price.}
State the question and choose a price the table accepts:
\begin{description}
  ℹtem[\textbf{Devil's Bargain}] Offer the GM a future complication to ``pay'' for deeper intel.
  ℹtem[\textbf{Accept SB}] The GM gains \textbf{+1 SB} now; you gain a strategic clue or reduced DV.
  ℹtem[\textbf{Spend Boon}] Spend \textbf{1 Boon} to lower DV by 1 on the research/social approach.
\end{description}

\paragraph{Creative Methods.}
Swap coin for color: poetry duels, shrine petitions, map‑reading at a parish stone. If your method sings with setting, take \textbf{+1 Effect} on the check.

\subsection{Timer Manipulation}
∈dex{timers}∈dex{pacing}
Players can nudge tension without rewriting stakes.
\begin{itemize}
  ℹtem \textbf{Slow} a visible Timer by \textbf{1 segment} by either \textbf{spending 1 Boon} \emph{or} \textbf{accepting +1 SB} into the scene (fiction must justify delay).
  ℹtem \textbf{Hasten} a visible Timer by \textbf{1 segment} by \textbf{spending 1 Momentum} (see \ref{subsec:momentum−banking}) \emph{or} \textbf{taking a Devil's Bargain}.
  ℹtem \textbf{Limit:} Once per player per scene. If three or more players affect the same Timer in a scene, the GM gains \textbf{+1 SB}.
\end{itemize}

\subsection{Complication Bargaining}
∈dex{complications!player−asked}
Invite the kind of trouble you want to play.
\begin{itemize}
  ℹtem Name a \emph{type}: social, physical, mystery, or moral (e.g., Valora court etiquette; Ubral scree; Isoka whispers; oath vs. mercy).
  ℹtem The GM frames the complication accordingly and grants you \textbf{+1d} on your next roll \emph{within} that trouble, or reduces DV by 1 if your approach leans into the specified texture.
\end{itemize}

\section{Faction Awareness}
∈dex{factions}∈dex{attitudes}

\subsection{Loyalty Recognition}
∈dex{loyalty scale}
Track a simple ladder \((−3…+3)\): \textit{Enemy, Hostile, Unfriendly, Neutral, Friendly, Supportive, Ally}. When you act \emph{with} an institution's aims, mark a \textbf{tick} toward the next step (GM pacing). When you betray a stated value, drop one step immediately and the GM banks \textbf{+1 SB} for future headaches.

\subsection{Cross−Cultural Synergies}
∈dex{synergy!cross−cultural}
Spotting a neat cultural combo (Zakov pilots + Kahfagia signals; Aeler engineering + Ecktoria charters) grants \textbf{+1d} once per scene the synergy is actively used.

\section{Advanced Techniques}
∈dex{advanced techniques}

\subsection{Momentum Banking}łabel{subsec:momentum−banking}
∈dex{momentum}
When your team resolves a Timer \emph{early} (segments unspent), bank \textbf{1 Momentum} per unused segment (max \textbf{2/session}). Spend 1 Momentum to:
\begin{itemize}
  ℹtem Gain \textbf{+1d} on a future approach tied to that victory (\emph{lessons learned}); or
  ℹtem \emph{Telescope} a travel beat (skip a routine obstacle the same route would present); or
  ℹtem Trigger a \emph{Prepared Move}: declare a sensible minor setup you plausibly arranged off‑screen.
\end{itemize}

\subsection{Escalation Management}
∈dex{escalation}
Trade heat for shape.
\begin{itemize}
  ℹtem \textbf{De‑escalate (spend 1 Boon):} Downgrade a Major consequence to Minor with a plausible concession (quiet tolls paid, harsh words eaten).
  ℹtem \textbf{Redirect (accept +1 SB):} Shift pressure to a new venue or actor you name; the GM places that SB there as attention.
  ℹtem \textbf{Truce (spend 1 Momentum):} Freeze a faction's hostility for one scene if you can cite a shared value (oath, rite, charter clause).
\end{itemize}

\section{Enhanced Coordination}
∈dex{coordination}∈dex{teamwork}

Teamwork isn't just about helping−−−it's about creating decisive moments through coordinated action.

\subsection{Core Principles}
\begin{itemize}
  ℹtem \textbf{Fiction First:} Describe how you help before rolling dice.
  ℹtem \textbf{Clear Stakes:} State what help changes and what risk the helper accepts.
  ℹtem \textbf{One Spotlight:} Resolve one acting character's roll; fold assistance into that action.
  ℹtem \textbf{Visible Costs:} Story Beats, Fatigue, and consequences are tracked openly.
\end{itemize}

\subsection{Basic Assistance}
\begin{itemize}
  ℹtem \textbf{Declare Help:} An ally states a concrete contribution (tools, opening, cover).
  ℹtem \textbf{Benefit:} Acting character gains \textbf{+1d} (up to the table's assist cap).
  ℹtem \textbf{Limits:} One assistant per PC per exchange by default.
  ℹtem \textbf{Cost:} Assistant accepts any risk named by the GM (SB, Fatigue, collateral).
\end{itemize}

\subsection{Cooperative Talents}
Certain talents enhance group coordination:

\textbf{Inspire (3 XP)} Once per scene, spend 1 Boon to choose one:
\begin{itemize}
  ℹtem Bonded Ally: That ally gains +1 Boon and +1d on their next roll.
  ℹtem Self: Gain +1d on your next roll.
  ℹtem Rally: Each other PC in Near gains +1d on their next roll.
  ℹtem Tactical Coordination: All allies currently acting gain +1 Position.
\end{itemize}

\textbf{Spotter's Mark (3 XP)} Aim a target (1 action). Until end of scene, allies in Near may claim +1d or +1 Effect once vs. that target.

\textbf{Battle Cant (2 XP)} Once per scene, establish silent signals. On the next coordinated action where at least two PCs act, those PCs gain +1 Position.

\textbf{Shield Wall (4 XP)} If you and at least one ally each wield a shield and are adjacent: as a Defend action, grant +1d Defend to all in the Wall and convert the first incoming Harm to Fatigue.

\subsection{Stacking Limits}
To prevent overwhelming combinations, a PC may benefit from at most \textbf{two cooperative effects} on the same action. Choose which apply when multiple sources could stack.

\subsection{Timing and Sequencing}
\begin{itemize}
  ℹtem \textbf{Declare Order:} Players state intent in any order; resolve the acting roll first.
  ℹtem \textbf{Ready/Overwatch:} Hold an action with a clear trigger; resolve before the next exchange.
  ℹtem \textbf{Refresh Windows:} ``Once/scene'' effects reset at scene end.
\end{itemize}

\subsection{Working With Followers}
Followers can assist for up to +3d (or +4d with the Exceptional Coordination talent). They cannot receive PC‑only benefits unless a talent states otherwise.

\subsection{Magic Coordination}
\begin{itemize}
  ℹtem \textbf{Ritual Support:} Allies can assist Weave or Cast phases if the fiction allows.
  ℹtem \textbf{Rite Preparation:} Helpers can provide components or maintain focus.
  ℹtem \textbf{Imbuement Synergy:} Multiple characters can prepare for shared magical effort.
\end{itemize}

\subsection{Best Practices for Players}
\begin{enumerate}
  ℹtem \textbf{Be Specific:} ``I'll flank left to give you advantage'' is better than ``I help.''
  ℹtem \textbf{Accept Risk:} Good assistance often means putting yourself in harm's way.
  ℹtem \textbf{Think Creatively:} Tools, positioning, and timing matter more than raw dice.
  ℹtem \textbf{Support Bonds:} Use coordination to strengthen character relationships.
  ℹtem \textbf{Keep Momentum:} Quick, clear assistance keeps scenes moving.
\end{enumerate}

\emph{Remember:} The best coordination enhances the story and makes everyone's moments shine brighter together.

\section{Character Creation Enhancements}
∈dex{character creation!enhancements}∈dex{bonds}∈dex{complications}
Players may take up to \textbf{2 Bonds} (+2 XP total) and up to \textbf{2 Starting Complications} (+4 XP total) for a cap of \textbf{36XP}. Each unresolved starting Complication adds \textbf{+1 banked SB} to early scenes until cleared. Favor \emph{storyful} picks (clan honour, guild debt, patron notice) over pure math.

\section{War Mount (Asset)}
łabel{subsec:war−mount}
∈dex{Assets!War Mount}∈dex{Mounts}

\textbf{Type:} Standard Asset \\
\textbf{XP Cost:} 8 XP \\
\textbf{Scale:} Personal \\
\textbf{Maintenance:} Regular care, feed, rest

A trained combat mount−−−warhorse, dune‑strider, battle elk, ash lizard, or equivalent−−−conditioned for violence, chaos, and discipline.

\paragraph{Passive Benefits}
While mounted:
\begin{itemize}
  ℹtem Increase overland travel speed and endurance.
  ℹtem Ignore the first movement‑based complication per session (mud, crowd, loose terrain).
  ℹtem Mounted positioning may override terrain restrictions at GM discretion.
\end{itemize}

\paragraph{Mounted Combat}
If the rider has the \textbf{Cavalier} talent (\ref{ch:talents}):
\begin{itemize}
  ℹtem Melee attacks made while moving \textbf{Far → Near}: \textbf{+2d}
  ℹtem Ranged attacks made while moving \textbf{Near → Far}: \textbf{+2d}
\end{itemize}

\paragraph{Boon Activation}
Spend \textbf{1 Boon} to invoke the mount dramatically:
\begin{itemize}
  ℹtem Break through a hostile line or crowd
  ℹtem Escape pursuit across hostile terrain
  ℹtem Gain \textbf{Controlled Position} for a mounted action
  ℹtem Advance or suppress a movement‑related Timer 
\end{itemize}

\paragraph{Risk and Wear}
On a \emph{Partial} or \emph{Miss} involving mounted actions, the GM may:
\begin{itemize}
  ℹtem Mark \textbf{Mount Fatigue}
  ℹtem Damage tack, barding, or terrain access
  ℹtem Threaten or compromise the Asset
\end{itemize}

\paragraph{Mount Fatigue}
If Mount Fatigue reaches 3:
\begin{itemize}
  ℹtem The mount cannot be used in combat until rested
  ℹtem Further strain risks permanent injury or loss
\end{itemize}

\paragraph{Promoting a War Mount}
A War Mount may be promoted to a \textbf{Follower} if:
\begin{itemize}
  ℹtem It is named and narratively significant
  ℹtem The rider takes a dedicated Bond or Talent
\end{itemize}
As a Follower, the mount can take actions independently (trample, defend, flee), suffer Harm, and may generate Boons through sacrifice or risk.

\subsection*{Pike Discipline (Major Talent, 4 XP) \textnormal{[FORMATION] [CONTROL]}}
∈dex{Talents!Anti−Cavalry}

\textbf{Requirements:} Melee 2+, Discipline or Command 1+

You are trained to break mounted assaults through spacing, timing, and fearlessness.

\paragraph{Effect}
When you \textbf{Brace} against a mounted opponent:
\begin{itemize}
  ℹtem Cancel all \textbf{Cavalier} movement bonuses
  ℹtem Mounted attackers suffer \textbf{−1d} on the attack
\end{itemize}

\paragraph{On Success}
\begin{itemize}
  ℹtem Reduce incoming Harm by 1
  ℹtem You may force the rider to \textbf{Disengage} or risk mount injury
\end{itemize}

\paragraph{On Partial}
\begin{itemize}
  ℹtem Reduce Harm by 1
  ℹtem Mark 1 Fatigue or lose formation
\end{itemize}

\paragraph{Limits}
\begin{itemize}
  ℹtem Requires braced footing or prepared terrain
  ℹtem Ineffective if flanked, surprised, or outnumbered
\end{itemize}

\section{Optional Player‑Character Death}
łabel{sec:pc−death−optional}
∈dex{Death}∈dex{Last Stand}∈dex{Death Timer }∈dex{Patron's Claim}∈dex{Inheritance}∈dex{Resurrection}

By default, Fate's Edge treats death as rare and dramatic. Use any of the following modules (singly or in combination) to tune how lethal your table feels. All options respect \textbf{Story Beats (SB)} and \textbf{Obligation} as core currencies.

\subsection*{Baseline (Default)}
PCs do not die on ordinary failures. Instead, they suffer \textbf{Harm}, Conditions, lost opportunities, or narrative costs. Death only occurs when a rule below is in play or the table agrees a scene warrants it.

\subsection*{Option A −−− Severe Harm Death}
\textbf{Trigger.} When a PC would take a \textbf{third} instance of Severe Harm (or escalate past the top of your harm track), they instead face death.\\
\textbf{Stave It Off.} The player may avoid death by choosing one: (1) mark \textbf{2 SB} \emph{and} take a \textsc{Maimed} permanent Condition, (2) accept a \textbf{Patron's Claim} (see Option D), or (3) convert the blow into a \textbf{Last Stand} (Option C).

\subsection*{Option B −−− Death Timer }
\textbf{Trigger.} Catastrophic consequences (falls, crushes, mortal wounds) fill a named \textbf{Death Timer } (4 or 6 ticks).\\
\textbf{While Ticking.} Actions that stabilise reduce the timer; taking further punishment advances it.\\
\textbf{When Full.} The character dies unless one of the following occurs immediately: spend \textbf{2 SB} to hold at full (buy a single action), accept a \textbf{Patron's Claim}, or another PC succeeds at \emph{Pull From the Brink} (risky, effect = timer −−2 on success).

\subsection*{Option C −−− Last Stand}
\textbf{Trigger.} On lethal harm or a full Death Timer , the player may declare a \textbf{Last Stand}.\\
\textbf{Effect.} For the remainder of the scene, the PC acts with \textbf{+1 Effect} and ignores new Harm. Each action automatically creates \textbf{1 SB}. When the scene ends, the character \textbf{dies} unless a \textbf{miracle} is secured (Patron rite, relic, or equivalent).

\subsection*{Option D −−− Patron's Claim}
\textbf{Trigger.} On death, the PC's Patron (or a circling power) intervenes.\\
\textbf{Bargain.} The GM offers 1−−3 non‑negotiable terms (e.g., \emph{Obligation +2}, \emph{forfeit a Gift}, \emph{become a vessel for a season}). If accepted, the PC lives; mark the costs immediately and record the \textbf{Claim} as an ongoing front. Refusal means the death proceeds.\\
\textbf{Note.} Claims should change the campaign; use sparingly.

\subsection*{Option E −−− Dramatic Exit & Inheritance}
\textbf{Dramatic Exit.} The player may choose a meaningful death that resolves a question or saves others.\\
\textbf{Inheritance.} The next PC created by that player inherits one of: (a) a \textbf{Relationship} (bond, contact, or rival), (b) a \textbf{Tool} (asset degraded one step), or (c) a \textbf{Lesson} (start with +1 XP toward a Talent used in the exit). Record how the world remembers them.

\subsection*{Option F −−− Return From Beyond}
\textbf{Trigger.} A body, a name, and a path (rite, gate, bargain).\\
\textbf{Cost.} Treat as a High Rite with Obligation appropriate to the transgression. Set DV using the Rites system. On success, the PC returns \textbf{changed}: apply a \textsc{Scar} (permanent Condition) and \textbf{1 SB} to the rescuer. On failure, choose: lose the body, or return with a Patron's Claim.

\begin{tcolorbox}[colback=gray!5!white,colframe=black,title=Quick Picks]
\small \textbf{Low lethality:} Baseline + Death Timer only.  \textbf{Heroic tragedy:} Severe Harm Death + Last Stand + Inheritance.  \textbf{Dark bargains:} Any combo with Patron's Claim.
\end{tcolorbox}

\section{Implementation Timeline}
∈dex{onboarding}∈dex{session flow}

\subsection*{Sessions 1−−3}
Investment Tracker, Information Trading (via Boons/Bargains), light Complication Bargaining.

\subsection*{Sessions 4−−6}
Timer Manipulation (Boon/Momentum/SB), Cultural Immersion, Faction Awareness, Momentum Banking.

\subsection*{Sessions 7+}
Cross‑Cultural Synergies, Escalation Management, advanced pacing choreography.

\begin{tcolorbox}[colback=blue!5!white,colframe=blue!75!black,title=At‑Table Prompts,fonttitle=\bfseries]
\textbf{Seed the Scene.} ``Ask for the kind of complication you want.''\\
\textbf{Name the Stakes.} ``Is this timer worth slowing? Who pays−−−Boon, SB, or Momentum?''\\
\textbf{Pay with Color.} ``What custom or rite do you invoke to make this work?''\\
\textbf{Close the Loop.} ``Mark your Investment; one sentence of what you learned.''
\end{tcolorbox}

\section{Downtime}
∈dex{downtime}∈dex{bookkeeping}

Between sessions, tend your garden−−−quietly moving the world.

\subsection{Bookkeeping}
\begin{itemize}
  ℹtem Allocate XP (respecting training days).
  ℹtem Update Assets/Followers (status: Maintained, Neglected, Compromised).
  ℹtem Track Boons (conversion cap remains 2 XP/session).
  ℹtem Note SB Debt from unresolved Complications.
\end{itemize}

\subsection{Activities}
\begin{itemize}
  ℹtem \textbf{Recovery:} Clear Harm/Exposure with scenes that show the work.
  ℹtem \textbf{Training:} Buy advances; narrate mentors, gyms, scriptoria.
  ℹtem \textbf{Research:} Lower DVs with good sources; log new leads.
  ℹtem \textbf{Social:} Strengthen Bonds; tune faction attitudes.
  ℹtem \textbf{Preparation:} Cache gear, sow rumours, line up ferries.
\end{itemize}

\subsection{Strategic Considerations}
Clear SB‑debt complications first; synchronise travel plans; decide which Asset gets love this interval, and which follower needs face‑time to avoid Neglect.

\section{Between‑Sessions Activities Log}
∈dex{logs!between sessions}∈dex{play aids}

Use or print the following trackers.

\subsection{Character Advancement}
⫕section*{XP Allocation}
\begin{itemize}
  ℹtem \textbf{Total XP Available}: _____
  ℹtem \textbf{Attributes Spent}: _____ (=___ days)
  ℹtem \textbf{Skills Spent}: _____ (=___ days)
  ℹtem \textbf{Remaining XP}: _____
\end{itemize}

⫕section*{Attribute Improvements}
\begin{tabular}{|l|c|c|c|c|}
ℎline
\textbf{Attribute} & \textbf{Old} & \textbf{New} & \textbf{Cost} & \textbf{Days} \\
ℎline
Body & & & & \\
Spirit & & & & \\
Presence & & & & \\
Wits & & & & \\
ℎline
\end{tabular}

⫕section*{Skill Improvements}
\begin{tabular}{|l|c|c|c|c|}
ℎline
\textbf{Skill} & \textbf{Old} & \textbf{New} & \textbf{Cost} & \textbf{Days} \\
ℎline
Arcana & & & & \\
Combat & & & & \\
Investigate & & & & \\
Lore & & & & \\
Move & & & & \\
Notice & & & & \\
Physique & & & & \\
Resolve & & & & \\
Stealth & & & & \\
Sway & & & & \\
Survival & & & & \\
ℎline
\end{tabular}

\subsection{Asset & Follower Management}
⫕section*{Assets}
\begin{tabular}{|l|c|c|c|}
ℎline
\textbf{Name} & \textbf{Tier} & \textbf{Status} & \textbf{Notes} \\
ℎline
 & & & \\
ℎline
\end{tabular}

⫕section*{Followers}
\begin{tabular}{|l|c|c|c|c|}
ℎline
\textbf{Name} & \textbf{Role} & \textbf{Harm} & \textbf{Exposure} & \textbf{Status} \\
ℎline
 & & & & \\
ℎline
\end{tabular}

\subsection{Bonds & Complications}
⫕section*{Bonds Updated}
\begin{itemize}
  ℹtem \textbf{With:} _____ −−− \textbf{Change:} _____
\end{itemize}
⫕section*{Complications}
\begin{tabular}{|p{5cm}|p{5cm}|}
ℎline
\textbf{Complication} & \textbf{Resolution or Status} \\
ℎline
 &  \\
ℎline
\end{tabular}

\subsection{Boons & Momentum Summary}
\begin{itemize}
  ℹtem \textbf{Boons Held:} __/5   \textbf{Converted to XP}: __ (max 2 XP/session)
  ℹtem \textbf{Momentum Banked:} __ (max 2/session)
\end{itemize}

\subsection{Campaign Timers}
\begin{itemize}
  ℹtem \textbf{Mandate:} __ / 6 \qquad \textbf{Crisis:} __ / 6
  ℹtem \textbf{Notables:} ________________
\end{itemize}